\documentclass[11pt,a4paper]{article}
\usepackage[margin=25mm]{geometry}
\usepackage{amsmath,amssymb,mathtools}
\usepackage{siunitx}
\usepackage{booktabs}
\usepackage{xcolor}
\usepackage[hidelinks]{hyperref}
\DeclareSIUnit\year{yr}
\definecolor{SolarAmber}{HTML}{E69F00}
\definecolor{CandidateTeal}{HTML}{009E73}
\definecolor{BaselineBlue}{HTML}{0072B2}
\title{Topology and Geometry Optimisation of Three-Dimensional Photovoltaic Canopies for Land-Constrained Singapore}
\author{Research project reconstruction}
\date{17 September 2026}
\begin{document}
\maketitle

\begin{abstract}
This project investigates three-dimensional photovoltaic geometries originating from a rotating mushroom-head concept for land-constrained Singapore. The principal hypothesis is spatial packing: additional active PV area may be placed above a constrained horizontal footprint while attempting to retain irradiation quality, sky visibility, ventilation and usable space beneath. No candidate geometry is assumed optimal. Exploratory calculations are distinguished from validated yield predictions.
\end{abstract}

\section{Solar-position equations}
For preliminary transparent calculations, the Cooper (1969) declination approximation is
\begin{equation}
\delta(n)=23.45^{\circ}\sin\!\left[\frac{360^{\circ}}{365}(284+n)\right].
\label{eq:cooper}
\end{equation}
Here $\delta$ is solar declination (degrees); $n$ is ordinal day of a non-leap year (dimensionless day index); $23.45^{\circ}$ is the approximate declination amplitude in the Cooper engineering correlation; $360^{\circ}$ is one complete angular cycle; $365$ is the non-leap-year period assumed by the approximation; and $284$ is its phase-offset constant. Equation~\eqref{eq:cooper} is an astronomical approximation, not an exact identity or a project-fitted model. Validated simulations should use a traceable higher-accuracy solar-position method.

Using local apparent solar time,
\begin{equation}
H=15^{\circ}\,\si{\per\hour}\left(t_{\mathrm{solar}}-12\,\si{\hour}\right).
\label{eq:hour-angle}
\end{equation}
Here $H$ is solar hour angle (degrees); $t_{\mathrm{solar}}$ is local apparent solar time (hours); $15^{\circ}\,\si{\per\hour}$ follows from $360^{\circ}/24\,\si{\hour}$; and $12\,\si{\hour}$ denotes apparent solar noon, where $H=0^{\circ}$. Civil clock time must not be substituted directly without longitude and equation-of-time correction.

Solar elevation satisfies
\begin{equation}
\sin\alpha=\sin\phi\sin\delta+\cos\phi\cos\delta\cos H.
\label{eq:elevation}
\end{equation}
Here $\alpha$ is solar elevation above the local horizon; $\phi$ is geographic latitude; $\delta$ is solar declination; and $H$ is solar hour angle. All angular quantities must use a consistent convention; the numerical implementation converts angles to radians before evaluating trigonometric functions.

\section{Direct and diffuse incidence}
For an infinitesimal PV element,
\begin{equation}
\mathrm dP_{\mathrm{dir}}=\eta(T)DNI(t)V(t)[\mathbf n\cdot\mathbf s(t)]_{+}\,\mathrm dA.
\label{eq:direct}
\end{equation}
Here $\mathrm dP_{\mathrm{dir}}$ is incremental direct-derived electrical power (W); $\eta(T)$ is conversion efficiency (dimensionless); $DNI$ is direct normal irradiance (\si{\watt\per\square\metre}); $V$ is direct visibility (dimensionless); $\mathbf n$ and $\mathbf s$ are unit vectors (dimensionless); $[x]_+=\max(0,x)$ is the positive-part operator; and $\mathrm dA$ is active PV area (\si{\square\metre}). The unit check is $(\si{\watt\per\square\metre})(\si{\square\metre})=\si{\watt}$.

The first diffuse approximation is
\begin{equation}
\mathrm dP_{\mathrm{diff}}=\eta(T)DHI(t)\frac{1+\cos\beta}{2}\,\mathrm dA.
\label{eq:diffuse}
\end{equation}
Here $DHI$ is diffuse horizontal irradiance (\si{\watt\per\square\metre}); $\beta$ is surface tilt from horizontal; and $(1+\cos\beta)/2$ is the unobstructed isotropic-sky view factor. The factor $1/2$ follows from ideal hemispherical geometry and is not an empirical Singapore coefficient.

\section{Paraboloidal mushroom}
The founding analytical geometry is
\begin{equation}
z(r)=h\left(1-\frac{r^2}{R^2}\right),\qquad k=\frac{h}{R}.
\label{eq:surface}
\end{equation}
Here $z$ is surface height (m); $r$ is radial coordinate (m); $h$ is apex height (m); $R$ is footprint radius (m); and $k$ is dimensionless curvature/aspect ratio.

Differentiating Eq.~\eqref{eq:surface},
\begin{equation}
\frac{\mathrm dz}{\mathrm dr}=-\frac{2hr}{R^2}.
\end{equation}
The slope is dimensionless. The numerical factor 2 arises exactly from differentiating $r^2$.

The axisymmetric area is
\begin{equation}
A_{\mathrm{PV}}=2\pi\int_0^Rr\sqrt{1+\frac{4h^2r^2}{R^4}}\,\mathrm dr.
\end{equation}
Here $A_{\mathrm{PV}}$ is active curved area (\si{\square\metre}); $2\pi r\,\mathrm dr$ is the annular area contribution; and the square-root term is the dimensionless slope correction.

Exact integration gives
\begin{equation}
A_{\mathrm{PV}}=\frac{\pi R^2}{6k^2}\left[(1+4k^2)^{3/2}-1\right].
\end{equation}
The factors 4 and 6 arise from the derivative and exact integral, not empirical fitting.

With $A_{\mathrm{foot}}=\pi R^2$,
\begin{equation}
\boxed{\Pi=\frac{A_{\mathrm{PV}}}{A_{\mathrm{foot}}}=\frac{(1+4k^2)^{3/2}-1}{6k^2}}.
\label{eq:packing-paraboloid}
\end{equation}
Here $\Pi$ is the dimensionless PV packing ratio and $A_{\mathrm{foot}}$ is horizontal footprint area (\si{\square\metre}). The limiting case $k\to0$ gives $\Pi\to1$, recovering a flat disk.

\section{Land multiplication}
Define
\begin{align}
\Pi&=\frac{A_{\mathrm{PV}}}{A_{\mathrm{land}}},\\
\eta_{\mathrm{pack}}&=\frac{E_{3D}}{\Pi E_{\mathrm{base}}}.
\end{align}
Here $A_{\mathrm{land}}$ is constrained horizontal site area (\si{\square\metre}); $E_{3D}$ and $E_{\mathrm{base}}$ are annual electrical energies (\si{\kilo\watt\hour\per\year}); and $\eta_{\mathrm{pack}}$ is dimensionless packed-PV productivity relative to the baseline.

Therefore
\begin{equation}
\boxed{M_L=\Pi\eta_{\mathrm{pack}}}.
\end{equation}
where $M_L$ is the dimensionless land multiplication factor.

\section{Mechanics}
Angular momentum and rotational kinetic energy are
\begin{equation}
L=I\omega,\qquad E_k=\frac12I\omega^2.
\end{equation}
Here $L$ is angular momentum (\si{\kilogram\metre\squared\per\second}); $I$ is mass moment of inertia (\si{\kilogram\metre\squared}); $\omega$ is angular velocity (\si{\radian\per\second}); and $E_k$ is rotational kinetic energy (J). The factor $1/2$ is the exact rigid-body kinetic-energy coefficient.

A first torque balance is
\begin{equation}
\tau=I\ddot\theta+\tau_f+\tau_g+\tau_w.
\end{equation}
Here $\tau$ is actuator torque (N m); $\ddot\theta$ is angular acceleration (rad s$^{-2}$); and $\tau_f$, $\tau_g$, $\tau_w$ are frictional, gravitational and wind torques (N m).

The first drag approximation is
\begin{equation}
F_D=\frac12\rho C_DA_{\mathrm{proj}}v^2.
\end{equation}
Here $F_D$ is drag force (N); $\rho$ is air density (kg m$^{-3}$); $C_D$ is dimensionless drag coefficient; $A_{\mathrm{proj}}$ is projected area (m$^2$); and $v$ is relative wind speed (m s$^{-1}$). The $1/2$ is the standard dynamic-pressure coefficient.

\section{Canonical topology experiment}
The initial reproducible design experiment uses
\begin{equation}
A_{\mathrm{foot}}\le\SI{1}{\square\metre},\qquad
\sum_iA_i\le\SI{2}{\square\metre},\qquad
0\le z_i\le\SI{2}{\metre}.
\end{equation}
Here $A_{\mathrm{foot}}$ is footprint area; $A_i$ is facet area; and $z_i$ is facet elevation. The numerical values 1, 2 and 2 are engineering design assumptions selected to define a reproducible comparison. They are neither Singapore regulatory limits nor discovered optima and must be varied in sensitivity analysis.

\section{Status and next step}
The Markdown report and governing-equation reference have been retrofitted to the equation-by-equation definition standard. This LaTeX file mirrors that policy. Before validated results are published, the source must be compiled, equations visually inspected in the PDF, and solar-position calculations upgraded/validated against a traceable higher-accuracy implementation.

\end{document}
